\begin{soln}
    We can set $\theta$ to be any number between the biggest blue
    point and the smallest red point. In fact, we can set $\theta$
    to be \emph{exactly} the largest blue point, as then all blue
    points are $\leq \theta$ as desired.

    We can do this in linear time by looping through all of the
    data points, keeping track of the largest blue point seen
    so far:

    \inputminted{python}{\solutiondir/pp01a-enum.py}

    A more ``pythonic'' approach is to ditch the \python{range}
    and use \python{zip}. \python{zip} takes two (or more) lists
    and returns pairs by ``zipping'' the lists together, elementwise:

    \inputminted{python}{\solutiondir/pp01a.py}

    Of course, we could always use a comprehension. The below
    filters out the red points and returns the max of the blues:

    \inputminted{python}{\solutiondir/pp01a-max.py}

    You could make the above into a one-liner, but I like the
    fact that the intermediate result is stored in \python{blue_points};
    it makes the code more readable, in my opinion. Avoid trying to show
    off by writing one-liners. A lot of people can cram code into
    one line, but it takes thought and care to write clean, readable code.

    All of these solutions are essentially the same, and run in
    linear time.
\end{soln}
